\documentclass[10pt]{article}
\usepackage{precalculus-article}
\usepackage{precalculus-boxes}
\newcommand{\TallMath}[1]{$\displaystyle #1\rule[-1.4em]{0pt}{3.2em}$}

\begin{document}

\pageheader{Unit 4, Lesson 4.5}{Guided Notes}
\namedateperiod

\vspace{0.1in}

\begin{objectivebox}
By the end of this lesson, I will be able to:
\begin{enumerate}[leftmargin=1.5em, itemsep=3pt]
    \item \writeline
    \item \writeline
    \item \writeline
\end{enumerate}
\end{objectivebox}

\vspace{0.08in}

\begin{vocabbox}
\termblanklong{Implicit equation}
\termblanklong{Explicitly defined function}
\termblanklong{Co-variation}
\termblanklong{Rate of change ($\Delta y / \Delta x$)}
\end{vocabbox}

\vspace{0.08in}

\begin{hookbox}
\textit{``The equation $x^2 + y^2 = 25$ is not a function --- but it describes a circle.
Can we still talk about how $x$ and $y$ change together? Can we split it into functions?''}

\vspace{4pt}
Is $x^2 + y^2 = 25$ a function? Why or why not?

\writeline
\end{hookbox}

\vspace{0.08in}

\begin{notesbox}{Part 1 — Graphing an Implicit Equation}

To graph an equation in two variables, \blank{8cm}

\vspace{0.1in}
\textbf{Example: $x^2 + y^2 = 25$}

Complete the table. (Hint: substitute $x$ and solve for $y$.)

\vspace{4pt}
\renewcommand{\arraystretch}{1.5}
\begin{tabular}{c|c|c}
$x$ & $y$ (positive) & $y$ (negative) \\\hline
$-5$ & \blank{2.2cm} & \blank{2.2cm} \\
$-3$ & \blank{2.2cm} & \blank{2.2cm} \\
$0$  & \blank{2.2cm} & \blank{2.2cm} \\
$3$  & \blank{2.2cm} & \blank{2.2cm} \\
$5$  & \blank{2.2cm} & \blank{2.2cm} \\
\end{tabular}

\vspace{0.12in}
This equation describes \blank{2cm} explicit function(s):

\vspace{4pt}
$y = $ \blank{4cm} \quad (upper half) \qquad $y = $ \blank{4cm} \quad (lower half)

\vspace{0.1in}
Solving for one variable in an implicit equation gives \blank{8cm}.

\end{notesbox}

\vspace{0.08in}

\begin{notesbox}{Part 2 — Co-variation}

\textbf{Definition:} Co-variation describes \blank{8cm}

\vspace{8pt}
\textbf{Example:} Points $(3, 4)$ and $(4, 3)$ are both on the upper half of $x^2 + y^2 = 25$.

\[
\frac{\Delta y}{\Delta x} = \frac{\blank{1.5cm}}{\blank{1.5cm}} = \blank{1.5cm}
\]

Is this ratio positive or negative? \blank{2cm}

What does this mean about how $x$ and $y$ change together?

\writeline

\vspace{0.1in}
\textbf{Rule:}
\begin{itemize}[leftmargin=1.5em, itemsep=3pt]
    \item If $\Delta y/\Delta x$ is \textbf{positive}, the variables \blank{5cm}.
    \item If $\Delta y/\Delta x$ is \textbf{negative}, the variables \blank{5cm}.
    \item If $\Delta y/\Delta x = 0$, the function has \blank{5cm} behavior.
    \item If $\Delta y/\Delta x$ is undefined, the function has \blank{5cm} behavior.
\end{itemize}

\end{notesbox}

\vspace{0.08in}

\begin{notesbox}{Part 3 — Second Example: $xy = 4$}

Solve for $y$: \quad $y = $ \blank{3cm}

\vspace{6pt}
\renewcommand{\arraystretch}{1.4}
\begin{tabular}{c|c}
$x$ & $y$ \\\hline
$1$ & \blank{1.8cm} \\
$2$ & \blank{1.8cm} \\
$4$ & \blank{1.8cm} \\
\end{tabular}

\vspace{6pt}
Compute $\Delta y/\Delta x$ from $x = 2$ to $x = 4$:

\vspace{4pt}
$\dfrac{\Delta y}{\Delta x} = \dfrac{\blank{2cm}}{\blank{2cm}} = \blank{2cm}$

\vspace{6pt}
As $x$ increases, $y$ \blank{3cm} (increases / decreases). Co-variation is \blank{3cm} (positive / negative).

\end{notesbox}

\end{document}
